\section{Análisis profundo del Agent Workflow}

\subsection{Descubrimiento y comprensión}

Cada agent empieza con intent detection: mapea la entrada en lenguaje natural a un structured goal. El equipo predefine slot templates para los intents comunes. La práctica que presenta Burak es usar primero un small model para hacer classification, y después un large model para hacer planning.

\subsection{Planeación y selección de herramientas}

El planner puede ser tanto un chain-of-thought prompt como una learned policy. Burak mostró un hybrid approach: primero se usa la policy para generar un plan tree, y después se usa el LLM para desarrollarlo paso a paso. Lo importante es el tool selection module: tiene que considerar el current status de la herramienta (si está sana), el estimated latency y el user context.

\subsection{Ejecución y revisión}

El Executor es, en realidad, una tool orchestration layer. Además de llamar a las API, tiene que manejar retries, fallbacks y confirm-with-user. Burak enfatiza que, para las operaciones financieras, siempre hay que agregar un explicit confirmation step, y registrar la decisión en el audit log.

\subsection{Cierre e iteración}

La última etapa del Workflow es la síntesis y el logging: se genera un summary para el usuario, se actualiza la memory, se registran todas las interacciones en el observability pipeline, y se dispara el training pipeline según el failure tag.

\subsection{Resumen de la sección}

Descomponer el agent workflow en cuatro pasos (descubrimiento, planeación, ejecución y cierre) ayuda a establecer métricas, monitoreo y revisión en cada etapa.
